\begin{theorem}[折叠多重度的有理 Hilbert--识别定理]\label{thm:conclusion-fold-inversecode-hilbert-recognizable}
沿用有限窗逆码的最小确定化表示记号。设
\[
d_m(x_1\cdots x_m)=u^\top T_{x_1}\cdots T_{x_m}v,
\qquad
x_i\in\{0,1\},
\]
其中 \((T_0,T_1,u,v)\) 为对应的最小线性实现。定义二元加权 Hilbert 级数
\[
\mathcal H(\xi,\eta)
:=
\sum_{m\ge 0}\ \sum_{x\in X_m}
d_m(x)\,\xi^{m-|x|_1}\eta^{|x|_1}.
\]
则有严格有理闭式
\[
\boxed{\;
\mathcal H(\xi,\eta)
=
u^\top\bigl(I-\xi T_0-\eta T_1\bigr)^{-1}v.\;}
\]
特别地，折叠多重度的全部字级统计构成一个严格有理的可识别级数；并且该线性实现的最小维数，恰等于 follower ambiguity classes 的基数。
\leanverified{matrix\_geom\_sum\_mul}
\leanverified{matrix\_mul\_geom\_sum}
\leanverified{weighted\_matrix\_sum\_unit}
\leanverified{paper\_conclusion\_fold\_inversecode\_hilbert\_recognizable}
\end{theorem}
\begin{proof}
由逐词矩阵实现，
\[
d_m(x_1\cdots x_m)=u^\top T_{x_1}\cdots T_{x_m}v.
\]
因此对每个固定 \(m\)，对所有长度 \(m\) 的词求和可得
\[
\sum_{x\in X_m}d_m(x)\,\xi^{m-|x|_1}\eta^{|x|_1}
=
u^\top(\xi T_0+\eta T_1)^m v.
\]
于是
\[
\mathcal H(\xi,\eta)
=
\sum_{m\ge0}u^\top(\xi T_0+\eta T_1)^m v
=
u^\top\Bigl(\sum_{m\ge0}(\xi T_0+\eta T_1)^m\Bigr)v,
\]
在形式幂级数意义下即为
\[
u^\top\bigl(I-\xi T_0-\eta T_1\bigr)^{-1}v.
\]
最后，最小线性实现维数等于右歧义类的有限完备集合基数，这正是正文关于最小确定化表示的极小性结论。证毕。
\end{proof}
